\subsection{The elliptic action coordinate as an explicit cup-product edge}

\begin{proposition}[E.35 --- Proved: the remaining elliptic coordinate is the Ross cup-product edge]
\label{E-ross-edge}
Retain the actual rational cubic class $[\eta]$ of
E.19, its compact splitting $V_{\mathbb Q_3}=H\oplus W$,
and the primitive determinant vector $w=u\wedge v$.
Put $x=-X/Z$, $y=-Y/Z$, so that $x^3+y^3=1$, and
use the unit tangent for $q=Z/X$ at $O=[1:-1:0]$.
The two functions
\[
 f=1-x,\qquad h=y-1
 \tag{E.35.1}
\]
are units on $U$. Their geometric Kummer maps
$\ell_f,\ell_h:V_{\mathbb Q_3}\to\mathbb Q_3(1)$
vanish on $H$. Let
\[
 b=\ell_f\wedge\ell_h:
       \bigwedge^2W\longrightarrow\mathbb Q_3(2),
 \qquad
 J:H\otimes\bigwedge^2W\longrightarrow\operatorname{Lie}_3(V),
 \quad J(a\otimes(r\wedge s))=[a,[r,s]].
 \tag{E.35.2}
\]
The map $J$ is injective. The two-dimensional invariant
elliptic summand $D_{HW^2}$ of E.28 is contained in its
image, and contraction defines an actual
$G_{\mathbb Q}$-equivariant isomorphism
\[
 \Phi=(1_H\otimes b)J^{-1}:
       D_{HW^2}\xrightarrow{\sim}H(2).
 \tag{E.35.3}
\]
This fixes a normalization of the remaining elliptic
coordinate, without choosing a scalar generator of its
one-dimensional arithmetic cohomology group.

Here is its value as a specified cohomology class.
Write $\Gamma=P\rtimes C_2$ for the actual free
pro-three arithmetic quotient of E.16, and
$\mathcal E=N\rtimes\Gamma$ for the geometrically
pro-three extension with its chosen tangent.
For the Hochschild--Serre edge map with the convention
specified in the proof, put
\[
 \begin{aligned}
 \rho&=\kappa(f)\smile\kappa(h)
              \in H^2_{\rm cts}(\mathcal E,\mathbb Q_3(2)),\\
 \operatorname{edge}_H(\rho)&\in
 H^1_{\rm cts}\bigl(\Gamma,
                   \operatorname{Hom}(H,\mathbb Q_3(2))\bigr),\\
 \iota_w(\lambda)&=u\lambda(v)-v\lambda(u).
 \end{aligned}
 \tag{E.35.4}
\]
In the last line the determinant twist gives the
equivariant map
$\iota_w:\operatorname{Hom}(H,\mathbb Q_3(2))\to H(1)$.
In (E.35.5), $\Phi_*$ is induced by the explicitly
twisted map $\Phi(-1):D_{HW^2}(-1)\to H(1)$.
If $[\eta]_{\rm ell}$ denotes the invariant elliptic
component in E.29, then
\[
 \boxed{\quad
  \Phi_*[\eta]_{\rm ell}
      =-\iota_{w*}\operatorname{edge}_H(\rho)
       \quad}\qquad\text{in }H^1(\Gamma,H(1)).
 \tag{E.35.5}
\]
Inflation gives the same equality for the original
$G_{\mathbb Q,\{3,\infty\}}$ action. Since
$\kappa_3(-1)=0$, the cup product in (E.35.4) is
also $\kappa(1-x)\smile\kappa(1-y)$, the Kummer
cup attached to the Ross symbol.
Formula (E.35.5) computes an exact comparison class.
Its numerical three-adic value, its vanishing or
nonvanishing, and the integral torsion coordinates
remain unevaluated. Proposition~\ref{E-kummer-values}
below evaluates the two Kummer coordinates of E.29.
\end{proposition}
\begin{proof}
First we fix every comparison needed by the cup
calculation. Let $\zeta=\zeta_3$ and label the cusps
\[
 A_i=[0:-\zeta^i:1],\quad
 B_i=[-\zeta^i:0:1],\quad
 C_i=[-\zeta^i:1:0]\quad(0\le i<3).
\]
Here $O=C_0$. The tangent-line equations give
\[
 \operatorname{div}(f)=3B_0-C_0-C_1-C_2,
 \qquad
 \operatorname{div}(h)=3A_0-C_0-C_1-C_2.
 \tag{E.35.6}
\]
For example, $X+Z=0$ meets the cubic with
intersection multiplicity three at $B_0$, whereas
$Z=0$ meets its three cusps simply. This proves
the first formula; the second follows from
$Y+Z=0$. The corresponding Kummer character on
a cusp meridian is its displayed order of
vanishing. The boundary module is
$\mathbb Q_3(1)^{\{A_i,B_i,C_i\}}$ modulo the
diagonal sum. Both order vectors sum to zero,
so they define precisely the maps $\ell_f,\ell_h$
on that quotient. On the compact summand the
maps are zero: they are Galois-equivariant, while
at good Frobenius the eigenvalues on $H$ have
weight one and that on $\mathbb Q_3(1)$ has
weight two. Thus the actual linear maps vanish
on $H$, not only on a chosen coordinate subspace.

At $O$ the branch with $Y/X=-1$ has
$Y/X=-(1+q^3)^{1/3}$. Hence
\[
 f=q^{-1}(1+q),\qquad
 h=q^{-1}\bigl((1+q^3)^{1/3}-q\bigr).
 \tag{E.35.7}
\]
Both leading coefficients are one. Their
Kummer cocycles therefore vanish on the specified
tangential arithmetic section. They factor through
$\mathcal E$: on the geometric group their
prime-three Kummer characters factor through $N$,
and their values on the section are zero. This
also specifies their values on the entire
semidirect product. A different rational tangent
changes them by arithmetic constant Kummer
classes; the compact edge of these added cross
terms is zero because both geometric characters
vanish on $H$. The fixed tangent in the statement
is the one used in the following literal formulas.

Choose lifts of a basis of $V$ in the rational
class-two Malcev Lie algebra. In its exponential
coordinates a geometric element is $(v,z)$,
with $z\in\bigwedge^2V$, and multiplication is
\[
 (v,z)(v',z')=
       (v+v',z+z'+\tfrac12v\wedge v').
 \tag{E.35.8}
\]
The genuine arithmetic action has the form
\[
 g(v,z)=\bigl(A_gv,
          (\bigwedge^2A_g)z+B_g(v)\bigr),
 \qquad C_g=B_gA_g^{-1}.
 \tag{E.35.9}
\]
Composition of these actual actions gives
$C_{gg'}=C_g+g\cdot C_{g'}$, with the usual
action on $\operatorname{Hom}(V,\bigwedge^2V)$.
Thus $C$ is a continuous one-cocycle. Changing
the graded coordinate lifts changes its class by
a coboundary in this same module.

The class-three bracket calculation of E.28 now
gives a more precise formula. Divide
$g[u,v]-\chi(g)[u,v]$ by $\chi(g)$. Its
$HW^2$ component is
\[
 \eta_{HW^2}(g)=
 J\bigl(u\otimes C_g^{WW}(v)
            -v\otimes C_g^{WW}(u)\bigr),
 \tag{E.35.10}
\]
where $C_g^{WW}$ is the restriction to $H$
followed by projection to $\bigwedge^2W$.
Indeed the determinant of $A_g|H$ is $\chi(g)$;
expanding the two cross brackets in (E.28.3)
cancels precisely that determinant. An inner
coordinate change contributes $a\wedge H$
to the degree-two correction, whose $WW$
projection is zero. Formula (E.35.10) therefore
also respects the specified quotient by inner
path ambiguity. The map $J$ is injective because,
in the tensor-algebra realization of the free
Lie algebra, the words whose first letter is
in $H$ recover the independent tensor
$a\otimes(r\wedge s)$ from $[a,[r,s]]$.

We compute the edge including its sign. Extend
$b$ by zero on the other multidegrees of
$\bigwedge^2V$. Write $e=(n,g)$, $e'=(m,g')$
and let $(v_n,z_n),(v_m,z_m)$ be their Malcev
coordinates. The antisymmetrized cup cocycle is
\[
 \omega(e,e')=
       \tfrac12 b(v_n\wedge A_gv_m).
\]
It represents $\rho$: graded commutativity
identifies the original cup with its
antisymmetrization, or the symmetric part is
the coboundary of the negative one-half product
of the two Kummer one-cochains. For the
continuous one-cochain $c(e)=-b(z_n)$, the
inhomogeneous differential and (E.35.8)--(E.35.9)
give exactly
\[
 \delta c(e,e')
    =\omega(e,e')+bB_g(v_m),
 \qquad
 \omega-\delta c=-bB_g(v_m)
                         =-(bC_g)(A_gv_m).
 \tag{E.35.11}
\]
The Tate action is included: $g\,b(z)=
b((\bigwedge^2A_g)z)$. The expression in the
last line is the mixed two-cocycle associated
to the left-module one-cocycle $-bC_g$.
We use this as the sign convention for the
Hochschild--Serre edge.

The edge is defined on the actual $\rho$ and
is an isomorphism here. The free pro-three
groups $N$ and $P$ have cohomology zero in
degree at least two for finite three-primary
coefficients. Their finite generating sets
make the cohomology at finite levels finite,
so the inverse-limit Mittag--Leffler argument
and rationalization give the same vanishing
for the Tate lattices and rational modules in
use. Averaging over $C_2$ preserves it.
Consequently the only nonzero term of total
degree two in the Hochschild--Serre sequence
for $N\rtimes\Gamma$ is
$H^1(\Gamma,\operatorname{Hom}(V,\mathbb Q_3(2)))$;
its possible outgoing degree-two differential
has target $H^3(\Gamma,\mathbb Q_3(2))=0$.
This proves the claimed edge comparison, with
the compact projection given by restriction
to the actual summand $H$. Applying $\iota_w$
to (E.35.11) and comparing (E.35.10) proves
(E.35.5) before the invariant normalization.

We finish by proving that the normalization
is an isomorphism on the required invariant
summand. This is a finite exact calculation
on the actual Fermat automorphisms. Let
$\mathcal A=\mu_3^2\rtimes S_3$ act by
diagonal scalings followed by coordinate
permutations. It is the same order-$54$ group
as $E[3]\rtimes\mu_6$ in E.23. In a rational
CM basis of $H$ put
\[
 R=\begin{pmatrix}0&-1\\1&-1\end{pmatrix}.
\]
The automorphism with diagonal exponents
$(a,b,0)$ and permutation $p$ acts on $H$ by
$\operatorname{sgn}(p)R^{a+b}$. This follows
by pulling back the differential
$d(X/Z)/(Y/Z)^2$: its multiplier is
$\operatorname{sgn}(p)\zeta^{a+b}$, and the
conjugate eigenvalue supplies the other CM
line. On cusp homology it is the literal
permutation of the displayed projective
points, followed by passage to the diagonal
quotient.

Use the images of
$A_0,A_1,A_2,B_0,B_1,B_2,C_0,C_1$ as a basis
of $W$, with $C_2$ their negative sum. If $P_p$
is any of the nine-point permutation matrices,
its quotient action is $\pi P_p L$, where
$L$ includes the first eight coordinate vectors
and $\pi(v)_i=v_i-v_8$. In particular
$\pi P_p=(\pi P_pL)\pi$, which verifies the
module comparison explicitly.
Let $e_{\mathcal A}$ be the average of the
$54$ operators on $H\otimes\bigwedge^2W$.
Consider the two columns
\[
 d_0=e_{\mathcal A}(a_0\otimes(A_0\wedge B_0)),
 \qquad
 d_1=e_{\mathcal A}(a_0\otimes(A_0\wedge B_1)).
\]
Here $a_0$ is the first CM basis vector used
only for this rank calculation; its scale does
not alter the determinant's nonvanishing or
the intrinsic map (E.35.3). The two order vectors
in (E.35.6) give
\[
 (1_H\otimes b)(d_0,d_1)=
     \begin{pmatrix}-1/3&1/3\\0&1/3\end{pmatrix},
 \qquad \det=-1/9.
 \tag{E.35.12}
\]
For a short check of every entry, sum first
over the nine diagonal scalings. For each of
the coordinate permutations $012$ and $102$
the resulting two-column matrix is
$\left(\begin{smallmatrix}-9&9\\0&9\end{smallmatrix}\right)$.
For each of the other four permutations it
is zero. Dividing their sum by $54$ gives
(E.35.12). The program
\texttt{ross\_contraction.py} also records the
entire integral numerator of the projector,
its rational inverse on these two columns,
and the conjugation check. All use exact
integer and rational arithmetic.

Thus the invariant part of $J$'s image has
dimension at least two. The whole $HW^2$
invariant space has dimension two by E.23,
so it equals this part. Formula (E.35.12)
proves that its contraction is an isomorphism.
The maps $J$ and $b$ are Galois-equivariant,
and Galois permutes the elements of
$\mathcal A$, so its average commutes with
Galois. The class $[\eta]_{HW^2}$ lies in the
image of this average by E.23; $B_w=[V,w]$
has at least two compact factors, hence has
zero $HW^2$ projection. Averaging therefore
does not change that cohomology class. This
completes the comparison to exactly the
elliptic component in E.29.

For the final identification, $-1$ has
compatible rational $3^r$th roots, so its
three-adic Kummer class is zero and replacing
$h$ by $1-y$ leaves $\rho$ unchanged.
The symbol $\{1-x,1-y\}$ is Ross's specified
symbol; his Theorem~1 proves that its multiple
by six is tame on the proper cubic and has
nonzero real regulator \cite{Ross1994}.
We use the explicit Kummer cup in (E.35.4),
not that real-regulator conclusion, to define
the three-adic class in (E.35.5).
The searched syntomic formula
\cite[Theorem 4.11]{Asakura2023} assumes
$p>\max(N,M)$; with $p=N=M=3$ this reads
$3>3$ and supplies no value for the present
class. Proposition~\ref{E-kummer-values} below
evaluates the two rational Kummer coordinates;
the numerical value of this exact three-adic
cup edge and the integral torsion remain.
The anchor is \S51.16,
folios 733v--736v, scans 172--178, and
LM7-E.23(i).
\end{proof}
